\section{Banco de veinte enunciados}\label{sec:corpus}
En estas fichas reúno intención, escritura, contexto, lectura esperada y contraste.
Uso la lectura esperada como clave de estudio y dejo el cuaderno para anotar
mi propia interpretación. Los referentes numéricos son
locales a cada ficha. Los conceptos y propiedades siguen el vocabulario de referencia.
\par\textbf{Cómo lo estudio.} Primero reviso las fichas. En otra sesión uso el
cuaderno, que muestra las expresiones sin sus soluciones. Conservo E01--E10
como controles de compatibilidad con la versión anterior.
\clearpage
\KalaFicha{T01}\vspace{5mm}\KalaFicha{T02}\clearpage
\KalaFicha{T03}\vspace{5mm}\KalaFicha{T04}\clearpage
\KalaFicha{T05}\vspace{5mm}\KalaFicha{T06}\clearpage
\KalaFicha{T07}\vspace{5mm}\KalaFicha{T08}\clearpage
\KalaFicha{T09}\vspace{5mm}\KalaFicha{T10}\clearpage
\KalaFicha{T11}\vspace{5mm}\KalaFicha{T12}\clearpage
\KalaFicha{T13}\vspace{5mm}\KalaFicha{T14}\clearpage
\KalaFicha{T15}\vspace{5mm}\KalaFicha{T16}\clearpage
\KalaFicha{T17}\vspace{5mm}\KalaFicha{T18}\clearpage
\KalaFicha{T19}\vspace{5mm}\KalaFicha{T20}
